This paper presents an Agent-Based Model for Central Clearing, including the understudied client clearing tier. The modelling of central clearing with a client tier through an Agent-Based Model is the first of its kind. It is built upon two layers: an exchange layer where fundamental, momentum, and noise traders put in orders on a Limit Order Book calibrated to reproduce the stylised facts of returns we see in empirical data; and a clearing layer where a Central Counterparty novates the trades placed on exchange, collecting margins and collateral and administering a default waterfall. The standard model is based on historical E-mini price data across a calm and a stressed period, and is then applied to study margin procyclicality. Simulations from the stressed model highlight the importance for non-bank clearing members to have an external source of funding, the inability to liquidate their own position to raise liquidity is a vulnerability. An observation made from the study of margin procyclicality is that higher constant margins do not remove risk, but relocate it from the client tier to the member tier. Members hit their capital limits and freeze the provisioning of clearing services, and the clients holding positions that they cannot exit end up defaulting. This highlights the importance for clients to have more than one clearing connection. An alternative model is later presented with stochastic and synthetic stressed paths, offering a calibrated ensemble of varied scenarios that can be employed for CCP systemic risk testing.